\chapter*{R\'esum\'e}
\addcontentsline{toc}{chapter}{R\'esum\'e}

Au B\'enin, de plus en plus d'entreprises, \'ecoles, administrations et PME se dotent
d'un site web sans disposer des comp\'etences ou des moyens n\'ecessaires pour
v\'erifier si celui-ci respecte les standards de s\'ecurit\'e de base. Les outils existants
(Mozilla Observatory, SSL Labs) sont souvent en anglais, dispers\'es, et peu adapt\'es
\`a un utilisateur non-technicien. Dans le cadre de notre stage acad\'emique, nous avons
con\c{c}u et d\'evelopp\'e \textbf{CyberGuard B\'enin}, une plateforme web qui automatise
le diagnostic de s\'ecurit\'e d'un site : l'utilisateur saisit une URL, et le syst\`eme
ex\'ecute une s\'erie de v\'erifications non intrusives (HTTPS, certificat SSL, en-t\^etes
de s\'ecurit\'e, cookies, exposition d'informations serveur), puis traduit ces r\'esultats
en un score sur 100, un niveau de risque et des recommandations compr\'ehensibles.
Le projet a d\'ebut\'e par une clarification du besoin et des objectifs, suivie d'une
phase de mod\'elisation UML (acteurs, cas d'utilisation, diagrammes de classes,
d'objets et de s\'equence) puis de la conception de la base de donn\'ees. La
r\'ealisation technique repose sur une API Laravel (PHP) c\^ot\'e backend et une
application React c\^ot\'e frontend. Ce projet propose ainsi une solution concr\`ete
pour d\'emocratiser l'acc\`es \`a un premier niveau d'audit de s\'ecurit\'e pour les
propri\'etaires de sites web au B\'enin.

\bigskip
\noindent\textbf{Mots-cl\'es~:} s\'ecurit\'e web, audit automatis\'e, HTTPS, Laravel,
React, UML.
